\documentclass[11pt]{article}
\usepackage[margin=1in]{geometry}
\usepackage{amsmath,amssymb,amsthm,mathtools}
\usepackage[T1]{fontenc}
\usepackage{lmodern,microtype}
\setlength{\emergencystretch}{3em}
\usepackage[hidelinks]{hyperref}
\newtheorem{theorem}{Theorem}
\newtheorem{lemma}{Lemma}
\newcommand{\C}{\mathbb C}
\newcommand{\R}{\mathbb R}
\newcommand{\norm}[1]{\left\lVert#1\right\rVert}
\title{Conjecture 00000007160\\Hoffman--Wielandt spectral variation}
\author{C0ldSmi1e}
\date{October 12, 2026}
\begin{document}
\maketitle
\begin{abstract}
We formulate the stated spectral Lipschitz claim using eigenvalues counted with
algebraic multiplicity and optimal Euclidean matching. For two complex normal
matrices of the same finite size, the squared matching cost is at most the
squared Frobenius norm of their difference. The argument uses complete
orthonormal eigenbases, Parseval's identity and the finite Birkhoff theorem.
The accompanying Lean project proves the estimate for arbitrary complete
eigenvalue lists, including repeated roots and zero dimension. Independent
source review, a fresh project build and a dependency/axiom audit passed.
These are local validation results, separate from maintainer acceptance.
\end{abstract}

\section{Statement and interpretation}
The original English conjecture is: ``The spectra of normal matrices are
Lipschitz with the Frobenius norm as constant.'' Both the English and Chinese
versions explicitly name Hoffman--Wielandt. The prose does not specify a
spectral metric. We use the conventional finite-dimensional formulation below.

For $A\in\C^{n\times n}$, let $A^*$ denote its conjugate transpose. Normality
means $AA^*=A^*A$. Set
\[
 \norm{A}_F^2=\sum_{i=1}^n\sum_{j=1}^n |A_{ij}|^2.
\]
A complete eigenvalue list $\alpha=(\alpha_1,\ldots,\alpha_n)$ records
\emph{algebraic multiplicity}, equivalently
$\det(tI-A)=\prod_i(t-\alpha_i)$. For two such lists define
\[
 d_2(\alpha,\beta)
 =\min_{\sigma\in S_n}
   \left(\sum_{i=1}^n |\alpha_i-\beta_{\sigma(i)}|^2\right)^{1/2}.
\]
The finite permutation set is nonempty, including $n=0$. Relabeling either
list simply permutes the set of matching costs, so this quantity depends on
the eigenvalue multisets. We interpret the Lipschitz claim as the following
pairwise inequality; no globally consistent labeling is asserted.

\begin{theorem}[Hoffman--Wielandt]
For every $n\geq0$, if $A,B\in\C^{n\times n}$ are normal, and $\alpha,\beta$
are their complete eigenvalue lists, then a permutation $\sigma\in S_n$ satisfies
\[
 \sum_i |\alpha_i-\beta_{\sigma(i)}|^2
 \leq \sum_{i,j}|A_{ij}-B_{ij}|^2.
\]
Consequently $d_2(\alpha,\beta)\leq\norm{A-B}_F$.
\end{theorem}
The Lipschitz constant is $1$ relative to the Frobenius norm. The theorem
allows repeated and nonreal eigenvalues, singular matrices, and pairs that
do not commute. Complex matrices also include real normal matrices considered
over $\C$. For $n=0$ all sums vanish. For $n=1$ equality holds, so the
constant cannot be lowered in positive dimension.

\section{Complete spectral data from normality}
Let $T$ be a normal linear operator on a finite-dimensional complex inner
product space. Its real and imaginary self-adjoint parts are
\[
 H=\frac{T+T^*}{2},\qquad K=\frac{T-T^*}{2i},\qquad T=H+iK.
\]
Normality gives $HK=KH$. The orthogonal eigenspace decomposition for $H$
is preserved by $K$: if $Hx=ax$, then $H(Kx)=K(Hx)=aKx$.
Diagonalizing the self-adjoint restriction of $K$ on each summand therefore
gives a complete orthonormal basis of joint eigenvectors. On a joint
eigenvector with eigenvalues $a,b$, the operator $T$ has eigenvalue $a+ib$.
This also covers the zero-dimensional space with its empty basis.

In this basis the matrix of $T$ is diagonal. Invariance of the characteristic
polynomial under a change of basis yields
\[
 \chi_T(t)=\prod_i(t-\alpha_i).
\]
Thus the resulting list is complete with algebraic multiplicity, not merely
a collection of values having eigenvectors. Any other complete list gives
the same multiset of roots, and hence differs only by a permutation.

\section{The squared Frobenius identity}
Use the inner product convention that is conjugate-linear in the first
argument and linear in the second. Let $(e_i)$ be a complete orthonormal
eigenbasis for $A$ with eigenvalues $\alpha_i$, and $(f_j)$ one for $B$ with
eigenvalues $\beta_j$. Define
\[
 w_{ij}=|\langle e_i,f_j\rangle|^2.
\]
Parseval's identity applied to each unit vector shows
\[
 w_{ij}\geq0,\qquad \sum_jw_{ij}=1,\qquad \sum_iw_{ij}=1.
\]
Thus $W=(w_{ij})$ is doubly stochastic. Since coordinates in the $e$ basis
diagonalize $A$, for every vector $x$,
$\langle e_i,Ax\rangle=\alpha_i\langle e_i,x\rangle$.
It follows that
\[
 \langle e_i,(A-B)f_j\rangle
 =(\alpha_i-\beta_j)\langle e_i,f_j\rangle.
\]
Apply Parseval in the $e$ basis and sum over $j$ to obtain
\begin{equation}\label{eq:weighted}
 \sum_{i,j}w_{ij}|\alpha_i-\beta_j|^2
 =\sum_j\norm{(A-B)f_j}^2
 =\norm{A-B}_F^2.
\end{equation}
For completeness, the last equality is basis invariance of this sum of
squared norms. If $D=A-B$ and $(u_k)$ is the standard orthonormal basis, then
\[
 \begin{aligned}
 \sum_j\norm{Df_j}^2
 &=\sum_k\sum_j|\langle u_k,Df_j\rangle|^2\\
 &=\sum_k\sum_j|\langle D^*u_k,f_j\rangle|^2
 =\sum_k\norm{D^*u_k}^2
 =\sum_{k,l}|D_{kl}|^2.
 \end{aligned}
\]
Every sum here is finite; exchanging its order needs no convergence argument.

\section{Selecting a matching}
The finite Birkhoff theorem writes a doubly stochastic matrix as a convex
combination of permutation matrices:
\[
 w_{ij}=\sum_{\sigma\in S_n}q_\sigma\,
       \mathbf1_{\{j=\sigma(i)\}},
 \qquad q_\sigma\geq0,\qquad\sum_\sigma q_\sigma=1.
\]
For any real cost matrix $C$, distributing finite sums gives
\[
 \sum_{i,j}w_{ij}C_{ij}
 =\sum_\sigma q_\sigma\sum_i C_{i,\sigma(i)}.
\]
The minimum of the finitely many permutation costs is at most this convex
average. Choose a minimizing permutation and set
$C_{ij}=|\alpha_i-\beta_j|^2$. Equation~\eqref{eq:weighted} gives the squared
inequality in the theorem. All quantities are nonnegative, so taking square
roots proves its final assertion. Relabeling the two spectral lists preserves
the set of permutation costs, giving the conclusion for arbitrary complete
lists.

\section{Formal correspondence and verification}
The project pins Lean 4.19.0 and Mathlib v4.19.0, at commit
\begin{center}\small\texttt{c44e0c8ee63ca166450922a373c7409c5d26b00b}\end{center}
{\raggedright Its public capstones,
in namespace \texttt{HoffmanWielandt}, are
\texttt{hoffman\_wielandt\_sq} and \texttt{hoffman\_wielandt}. They accept
arbitrary matrices, their explicit normality equations and arbitrary complete
characteristic-polynomial eigenvalue lists. Separate existence theorems
construct such lists from normality, so the formulation is not vacuous.\par}

\begin{itemize}\raggedright
\item \texttt{NormalSpectral.lean} proves the complete normal eigenbasis,
diagonal coordinates and characteristic-polynomial factorization.
\item \texttt{Frobenius.lean} proves the Parseval identities, doubly stochastic
overlap property and exact weighted Frobenius identity. Its explicit norm is
also proved equal to Mathlib's dedicated Frobenius norm instance.
\item \texttt{Matching.lean} proves the permutation-cost bound from Birkhoff's
theorem. \texttt{Reindexing.lean} proves that equal products of linear factors
give a permutation of the complete lists, using root multiplicities.
\item \texttt{HoffmanWielandt.lean} transfers between matrices and operators,
assembles the estimate and handles arbitrary eigenvalue enumerations.
\texttt{Verification.lean} checks the expanded statement, empty dimension,
repeated nonreal roots, reversed enumeration, signed costs, sharp scalar
equality and explicitly non-Hermitian, noncommuting normal matrices.
\end{itemize}

A fresh private project copy passed \texttt{lake --no-cache --wfail build}.
All six mathematical source files and the Lake configuration were then
replayed directly with warnings treated as errors. Both theorem and
verification libraries are default build targets. The audit inventoried
49 originating-module declarations and a dependency graph of 30,338
declarations, and bound 4,152 imported module identities to the recorded
runtime, pinned dependency copies and rebuilt local modules. All 36 safe originating declarations,
including the four main capstones, use only \texttt{propext},
\texttt{Classical.choice} and \texttt{Quot.sound}, with no unsafe proof
dependency. There are no admitted proofs, custom axioms or
\texttt{native\_decide} calls in the submitted mathematical sources.

The remaining 13 inventoried definitions are compiler-generated runtime
helpers for the two explicit matrices in \texttt{Verification.lean}. Their
exact names and dependency information were inspected separately. No safe
originating declaration or capstone reaches them, and they are excluded from
mathematical evidence. The raw pending-classification audit record is
preserved; the separate final classification completes it. Generic negative
controls confirmed rejection of admitted proofs, direct and indirect custom
axioms, native-computation proof shortcuts and unsafe dependencies.

All nine pinned dependency source trees were checked against their Git
objects. This review used recorded stock dependency caches; it did not rebuild
all of Mathlib or bootstrap the supplied Lean runtime. The dependency inspector
and graph-consistency checker are not independent kernel recheckers: the
fresh source compilations supply Lean's kernel checks. Independent mathematical
review read the entire final source and the report's argument against both
original language versions.

The separate standard-library Python checker uses exact rational complex
arithmetic. It checks 361 generated normal pairs in dimensions $0$--$3$,
verifies their complete characteristic polynomials by the determinant formula,
and enumerates every matching. Among these pairs, 194 do not commute, 335
have a non-Hermitian member, 127 have repeated eigenvalues and 188 satisfy a
strict inequality; these categories overlap. Four additional semantic controls
exercise enumeration, multiplicity, the norm choice and the normality
hypothesis. A separate review inspected the entire checker, replayed it and
passed 71 implementation probes. These are finite sanity checks and do not
establish the theorem for arbitrary matrices. No numerical example is a
premise in the general Lean proof.

\section{Reproduction and records}
With Lean 4.19.0 and the locked dependency caches available, run these commands
from the personal submission directory:
\begin{verbatim}
lake --wfail build
python3 -B auxiliary/check_examples.py
\end{verbatim}
Compile the report source with LaTeX or Tectonic.

The \texttt{verification} directory contains source identities, independent
reviews, the engineering summary and runtime-helper classification. Its
compressed raw-evidence archive retains the full dependency export, import
identities and build/replay receipts, without dependency caches. Included
verifier sources contain deliberately invalid negative-control fixtures;
these are test inputs, never parts of the mathematical proof. These records
document local validation, not maintainer acceptance.

\end{document}
